% !PRINTED 201
\weberpage{201}
We take any one of the equations (1)

\weq{8}{a = \Theta(a)}

and in the equation

\weq{9}{a_1 = \psi_1(a)}

replace $a$ by the equal value $\Theta(a)$. We thus obtain

\weq{10}{a_1 = \psi_1 \Theta(a).}

But by (9) we have $a = \psi_1^{-1}(a_1)$, and consequently from (10) there results:

\weq{11}{a_1 = \psi_1 \Theta \psi_1^{-1}(a_1).}

From this one sees, since we have $\nu - 1$ distinct substitutions $\Theta$, that $a_1$ is a $\nu$-fold pole, and indeed one belonging to the group $\psi_1 Q \psi_1^{-1}$ conjugate to $Q$. For this reason we call

\weq{12}{a, a_1, a_2, \ldots, a_{\mu-1}}

\gesperrt{a system of conjugate $\nu$-fold poles of $G$.}

An arbitrary substitution $\chi$ of $G$ can, by (6), always be brought into the form $\psi_i \Theta$, so that $\Theta$ belongs to $Q$, and from this it follows that $\chi(a) = \psi_i \Theta(a) = a_i$; and likewise, if one sets $\psi_i \Theta \psi_k = \psi_h \Theta'$, so that $\Theta'$ also belongs to $Q$:

\[ \chi(a_k) = \psi_i \Theta \psi_k(a) = \psi_h \Theta'(a) = a_h. \]

Since, further, two quantities $\chi(a_h), \chi(a_k)$ can be equal to one another only when $a_h = a_k$, the theorem follows:

5. \gesperrt{The two sequences}

\[ a, \quad a_1, \quad a_2, \quad \ldots, \quad a_{\mu-1} \]
\[ \chi(a), \quad \chi(a_1), \quad \chi(a_2), \quad \ldots, \quad \chi(a_{\mu-1}) \]

\gesperrt{agree with one another, apart from the order, whatever substitution from $G$ may be taken for $\chi$.}

If now $b$ is a pole not contained in the system (12), then $\chi(b)$ also cannot occur in (12), and from this it results that two systems of conjugate poles, if they are not entirely identical, have no pole in common.

Accordingly we can arrange all the poles of the group $G$ into systems of conjugate poles, and by \S.\,51 we obtain their total number $2n - 2$ if we count each $\nu$-fold pole $(\nu - 1)$ times.
